% Fragmento integrable. Preambulo anfitrion: amsmath, amssymb, longtable, hyperref.
\section{Identification analytique des centres initiaux}
\label{hmtfg:fragmento:identificacion-analitica-de-los-centros-originales}

\subsection{Réalisation de degré deux et identification métrique des centres}
\label{hmtfg:clasificacion:8}

L’égalité entre le premier paramètre de l’itinéraire infini et le croisement stable, démontrée par première sortie, fournit
\begin{equation}
\boxed{a_*:=\min\Lambda_\tau=\lambda_*.}
\label{hmtfg:clasificacion:eq:H}
\end{equation}

La famille, le profil dodécaphasique, les retours et le sélecteur restent exactement tels qu’ils sont dans les sections~\ref{hmtfg:clasificacion:1}--\ref{hmtfg:clasificacion:5}. Les résultats externes employés ci-dessous reconnaissent la dynamique de cet objet déjà construit. Leurs valeurs universelles n’interviennent pas dans la sélection des coefficients ni des paramètres de la famille HMT.

\textbf{Théorème de transfert métrique.} L’identification \eqref{hmtfg:clasificacion:eq:H}, jointe aux bornes d’entrée, d’admissibilité et de transversalité démontrées précédemment, implique que la suite initiale des premières racines nouvelles satisfait, pour certaines constantes \(C>0\) et \(0<\theta<1\),
\begin{equation}
\boxed{\lambda_*-\lambda_n=C\delta_{\mathrm F}^{-n}
\bigl(1+O(\theta^n)\bigr),}
\label{hmtfg:clasificacion:eq:14}
\end{equation}
où \(\delta_{\mathrm F}>1\) est la valeur propre instable de l’opérateur de doublement au point fixe reconnu. En particulier,
\begin{equation}
\frac{\lambda_{n-1}-\lambda_{n-2}}
     {\lambda_n-\lambda_{n-1}}
=\delta_{\mathrm F}+O(\theta^n).
\label{hmtfg:clasificacion:eq:15}
\end{equation}
La suite coïncide à partir d’un certain rang avec les centres de la continuation locale, une fois que les deux indices désignent la même période physique. La preuve conserve le sélecteur de première racine de la section~\ref{hmtfg:clasificacion:4}.

\subsubsection{Obtention d’un domaine commun de type quadratique}
\label{hmtfg:clasificacion:8.1}

Notons \(g\) le point fixe de la renormalisation et
\(\Gamma(\lambda)=R^4f_\lambda\). La carte analytique utilisée dans le certificat est
\[
f(x)=1-x^2V\!\left(\frac{x^2-1}{2.5}\right),\qquad
V(z)=\frac{u}{10}+\sum_{k\ge1}\nu_kz^k,
\]
avec la norme des différences \(|\Delta u|+\|\Delta\nu\|_1\). Les inégalités d’admissibilité prouvent que \(g\) est paire, analytique, unimodale, quadratique et de type doublement ; de plus, \(g(x)=\varphi(x^2)\) avec \(\varphi'<0\) dans \([0,1]\).

Le théorème 2.1 de de Faria–de Melo–Pinto, appliqué à l’ensemble de combinatoires constitué uniquement du doublement, fournit un ensemble limite à un seul point avec une extension de type quadratique. Puisque \(Rg=g\), la convergence réelle de ce théorème identifie ce point à \(g\). Le lemme 3.2 fournit, dans un voisinage analytique de \(g\), une itérée finie de renormalisation avec un domaine de type quadratique et un module uniformément positif. Les références précises sont les sections 2.1–2.3, p. 736–738, et la section 3.1, lemme 3.2, p. 745, de \href{https://annals.math.princeton.edu/wp-content/uploads/annals-v164-n3-p01.pdf}{de Faria–de Melo–Pinto, \emph{Global hyperbolicity of renormalization for \(C^r\) unimodal mappings}}.

Pour vérifier l’appartenance au voisinage de ce lemme, soit \(\Omega_\epsilon\) un voisinage complexe suffisamment petit de \([-1,1]\), avec
\[
\rho_\epsilon:=\sup_{x\in\Omega_\epsilon}
\left|\frac{x^2-1}{2.5}\right|<1,
\qquad M_\epsilon:=\sup_{x\in\Omega_\epsilon}|x|^2<\infty.
\]
La restriction de la carte certifiée satisfait
\begin{equation}
\|\Delta f\|_{\Omega_\epsilon}
\le M_\epsilon\left(\frac{|\Delta u|}{10}
          +\rho_\epsilon\|\Delta\nu\|_1\right)
\le C_\epsilon\|\Delta f\|_{\mathcal A}.
\label{hmtfg:clasificacion:eq:16}
\end{equation}
Le certificat donne
\[
\|R^m\Gamma(\lambda_*)-g\|_{\mathcal A}
<0.001666(0.83996)^m.
\]
On choisit donc un entier fixe \(m\) qui place cette image dans le voisinage du lemme 3.2. La continuité de la composition finie fournit un intervalle \(I_*\) autour de \(\lambda_*\) où les mêmes opérations sont définies. Si \(N\) est l’itérée fournie par ce lemme, nous posons
\begin{equation}
d=4+m+N,\qquad
\mathcal F_\lambda=R^df_\lambda,\quad \lambda\in I_*.
\label{hmtfg:clasificacion:eq:17}
\end{equation}
La famille \(\mathcal F\) possède des représentants de type quadratique de degré deux, sur un domaine commun et avec un module uniformément positif. Sa dépendance analytique réelle admet une complexification locale du paramètre. Le nombre \(d\) est fixe et compte les doublements effectués avant d’appliquer le résultat paramétrique.

Cette construction fixe un domaine commun et un nombre fini de retours \(d\), suffisants pour appliquer le théorème paramétrique.

\subsubsection{Transport de la transversalité certifiée}
\label{hmtfg:clasificacion:8.2}

La direction de \(\Gamma\) satisfait la condition de cône
\[
\|\dot\nu\|_1\le\tfrac14|\dot u|,\qquad \dot u\ne0.
\]
Les bornes différentielles certifiées impliquent l’invariance de ce cône et
\begin{equation}
|\dot u_{k+1}|\ge4.381|\dot u_k|
\label{hmtfg:clasificacion:eq:18}
\end{equation}
le long de l’orbite de \(\Gamma(\lambda_*)\). Tous ses points restent dans le domaine certifié. En particulier, la direction transportée par les étapes finies de \eqref{hmtfg:clasificacion:eq:17} conserve une composante expansive non nulle.

La classe hybride du point fixe coïncide avec sa variété stable dans l’espace des germes de type quadratique : \href{https://www.maths.tcd.ie/EMIS/journals/Annals/149_2/lyubich.pdf}{Lyubich, \emph{Feigenbaum–Coullet–Tresser universality and Milnor's hairiness conjecture}, théorème 6.1, p. 371–372}. La transversalité reste valable lors du passage à cet espace, comme on peut le vérifier sans identifier des normes de Banach différentes. En effet, si \(\partial_\lambda\mathcal F_{\lambda_*}\) était tangente à la classe hybride, elle serait la tangente d’un disque analytique contenu dans cette classe. La convergence exponentielle uniforme de la renormalisation sur un sous-disque compact, suivie de l’estimation de Cauchy en son paramètre, ferait converger vers zéro l’évaluation en \(x=1\) de ses tangentes renormalisées. Cependant, dans la carte précédente,
\[
\dot f(1)=-\dot u/10,
\]
et \eqref{hmtfg:clasificacion:eq:18} fait croître sa valeur absolue. Les deux dérivées représentent le même germe et ont la même évaluation réelle. La contradiction démontre
\begin{equation}
\partial_\lambda\mathcal F_{\lambda_*}
\notin T_{\mathcal F_{\lambda_*}}\mathcal H_{c_\infty}.
\label{hmtfg:clasificacion:eq:19}
\end{equation}
Ainsi, \(S=\{\mathcal F_\lambda\}\) est une transversale analytique complexe à la classe hybride de la limite de doublement.

\subsubsection{Identification à partir d’un certain rang des premières racines par leur période}
\label{hmtfg:clasificacion:8.3}

Notons \(c_r\) le centre quadratique canonique de période \(2^r\), obtenu par \(r\) doublements successifs du centre de période un. Ce symbole désigne le paramètre de la famille quadratique de reconnaissance ; il se distingue des fonctions orbitales \(c_j(\lambda)=f_\lambda^j(0)\) des sections précédentes. On écrit \(c_r^{\rm quad}\) pour conserver cette distinction.

Le théorème de continuation globale et \eqref{hmtfg:clasificacion:eq:H} donnent \(\lambda_n\to\lambda_*\), d’où \(\lambda_n\in I_*\) pour \(n\) suffisamment grand. La classification et les intervalles restrictifs de la section~\ref{hmtfg:clasificacion:2} montrent que, lorsque \(n>d\),
\begin{equation}
\mathcal F_{\lambda_n}=R^df_{\lambda_n}
\quad\text{a une période critique exacte }2^{n-d}
\quad\text{et itinéraire }W_{n-d}0.
\label{hmtfg:clasificacion:eq:20}
\end{equation}
Son orbite critique reste dans l’intervalle réel invariant, contenu dans le domaine de type quadratique ; son ensemble de Julia rempli est donc connexe. Le redressement conserve l’orbite postcritique et les branches réelles. En parcourant ses retours restrictifs, on obtient \(n-d\) doublements canoniques et, à la fin, un point critique fixe. Ce dernier se redresse au centre \(0\) ; les inverses successifs du redressement des copies de doublement donnent l’unique centre \(c_{n-d}^{\rm quad}\). Il s’agit de l’application concrète du paramétrage par combinatoire et de l’homéomorphisme de redressement des copies, décrits dans Lyubich, sections 5.1–5.3, p. 362–364. On conclut
\begin{equation}
\mathcal F_{\lambda_n}\in\mathcal H_{c_{n-d}^{\rm quad}}.
\label{hmtfg:clasificacion:eq:21}
\end{equation}

Le théorème 7.4 de Lyubich, p. 387–388, donne une unique intersection locale de \(S\) avec chaque classe \(\mathcal H_{c_r^{\rm quad}}\) pour \(r\) suffisamment grand. Ses hypothèses de bornes complexes a priori sont satisfaites pour la cascade quadratique réelle de doublement, conformément au théorème 5.6, p. 367, du même travail. Notons \(\zeta_r\) le paramètre de cette intersection. Les équations \eqref{hmtfg:clasificacion:eq:20}–\eqref{hmtfg:clasificacion:eq:21}, la convergence vers le croisement et l’unicité locale impliquent
\begin{equation}
\boxed{\lambda_n=\zeta_{n-d}\quad\text{pour tout }n\text{ suffisamment grand}.}
\label{hmtfg:clasificacion:eq:22}
\end{equation}
L’unicité utilisée ici concerne \textbf{chaque classe hybride canonique} dans un voisinage fixe du croisement. Elle ne se déduit ni de la seule existence d’une cascade locale ni de l’unicité du croisement avec la feuille stable. L’équation \eqref{hmtfg:clasificacion:eq:22} identifie le sélecteur initial, déjà défini globalement, aux intersections locales ; elle n’introduit pas de sélecteur de substitution.

En particulier, soit \(\mu_j\) la cascade locale de la section \ref{hmtfg:escalado:8}, et soit \(s\) l’entier fixe déterminé par sa période physique exacte :
\[
\operatorname{per}_{\rm crit}(f_{\mu_j})=2^{j+s}.
\]
Dans la présentation avec la courbe \(R^4f_\lambda\) et la cible de période deux, \(s=5\) ; les étapes fixes supplémentaires utilisées pour localiser la cible modifient \(s\). Par \eqref{hmtfg:clasificacion:eq:21} et la même unicité,
\begin{equation}
\mu_j=\zeta_{j+s-d},\qquad
\boxed{\lambda_n=\mu_{n-s}\quad(n\text{ suffisamment grand}).}
\label{hmtfg:clasificacion:eq:23}
\end{equation}
Le décalage \(d\) de la préparation de type quadratique et le décalage \(s\) des périodes de la cible sont des objets distincts. L’identification s’effectue par la période physique, et non en égalant ces deux entiers.

\subsubsection{Reste en puissance par holonomie transversale}
\label{hmtfg:clasificacion:8.4}

Le lemme 7.3 de Lyubich, p. 384, avec démonstration p. 384–387, apporte une régularité \(C^{1+\beta}\)-conforme, pour un certain \(\beta>0\), de l’holonomie entre transversales à la classe hybride du point de Feigenbaum. Sa définition p. 384 concerne les lieux de connexité sur les transversales ; ils contiennent notamment les centres considérés ici. Nous pouvons réduire l’exposant de sorte que \(0<\beta\le1\).

Soit \(h\) l’holonomie de \(S\) vers la variété instable et soit \(z\) une coordonnée analytique qui linéarise la restriction unidimensionnelle de la renormalisation :
\begin{equation}
z(g)=0,\qquad z(Rq)=\delta_{\mathrm F} z(q).
\label{hmtfg:clasificacion:eq:24}
\end{equation}
L’existence de cette coordonnée s’obtient en appliquant la linéarisation d’un germe holomorphe répulsif, de multiplicateur \(\delta_{\mathrm F}>1\), à la variété instable analytique. La direction expansive de \eqref{hmtfg:clasificacion:eq:18}, l’invariance du cône et le retour stationnaire identifient ce multiplicateur à la valeur propre instable positive employée dans le certificat.

Si \(q_r\) est le centre de la classe \(\mathcal H_{c_r^{\rm quad}}\) situé sur la variété instable, la compatibilité du redressement et de la renormalisation donne \(Rq_r=q_{r-1}\). Cette identité est celle utilisée dans la démonstration du théorème 7.4, p. 388. Par \eqref{hmtfg:clasificacion:eq:24}, il existe \(b\ne0\), qui absorbe tout indice initial fixe, tel que
\begin{equation}
z(q_r)=b\delta_{\mathrm F}^{-r}
\label{hmtfg:clasificacion:eq:25}
\end{equation}
pour tout \(r\) suffisamment grand.

La régularité de \(h\) et la transversalité \eqref{hmtfg:clasificacion:eq:19} donnent, sur le lieu de connexité proche de \(\lambda_*\),
\begin{equation}
H(\lambda):=z\bigl(h(\mathcal F_\lambda)\bigr)
=A(\lambda-\lambda_*)
+O\!\left(|\lambda-\lambda_*|^{1+\beta}\right),
\qquad A\ne0.
\label{hmtfg:clasificacion:eq:26}
\end{equation}
Puisque \(H(\zeta_r)=b\delta_{\mathrm F}^{-r}\), l’estimation inférieure
\(|H(\lambda)|\ge |A|\,|\lambda-\lambda_*|/2\) dans un voisinage suffisamment petit implique d’abord
\(|\zeta_r-\lambda_*|=O(\delta_{\mathrm F}^{-r})\). En substituant cette borne dans le reste de \eqref{hmtfg:clasificacion:eq:26}, on obtient
\begin{equation}
\zeta_r-\lambda_*=\frac bA\delta_{\mathrm F}^{-r}
+O\!\left(\delta_{\mathrm F}^{-(1+\beta)r}\right).
\label{hmtfg:clasificacion:eq:27}
\end{equation}
Nous appliquons \eqref{hmtfg:clasificacion:eq:22} avec \(r=n-d\). Le décalage fini modifie le coefficient principal et la constante du reste. Puisque \(\lambda_n<\lambda_*\), on obtient \eqref{hmtfg:clasificacion:eq:14}, avec
\[
C=-\frac bA\delta_{\mathrm F}^d>0,
\qquad \theta=\delta_{\mathrm F}^{-\beta}\in(0,1).
\]
Enfin,
\[
\lambda_n-\lambda_{n-1}
=C(\delta_{\mathrm F}-1)\delta_{\mathrm F}^{-n}\bigl(1+O(\theta^n)\bigr),
\]
ce qui démontre \eqref{hmtfg:clasificacion:eq:15}. \(\square\)


\subsubsection{Portée quantitative du théorème}
\label{hmtfg:clasificacion:8.5}

L’identification des limites démontrée par première sortie active le transfert métrique pour les premières racines globales du secteur uniforme. Le taux opératoriel \(0.83996\) contrôle les applications renormalisées sur la feuille stable. Le taux paramétrique \(\theta=\delta_{\mathrm F}^{-\beta}\) découle de la régularité de l’holonomie transversale. Les deux estimations conservent leur variable et leur domaine.

Le théorème établit l’existence de \(\beta,\theta,C\), du voisinage \(I_*\), du nombre fini de retours \(d\) et d’un indice \(n_0\) à partir duquel les centres coïncident. Leurs valeurs numériques individuelles ne sont pas attribuées dans cet énoncé asymptotique. Cette formulation complète la limite et son reste en puissance ; l’évaluation numérique de ces témoins constitue une spécialisation quantitative distincte.
